\chapter[Fisher Information Geometry]{Relation with Fisher Information Geometry}
Throughout this chapter let let $(\mathcal P,g)$ denote the posterior manifold equipped with the pullback Riemannian metric introduced in Chapter 7.
The purpose of this chapter is to compare the geometric structure
induced by Weighted Incidence Structures with the classical Fisher
information geometry.
\begin{remark}
The results of this chapter establish structural
relationships. They do not assert that the pullback metric always
coincides with the Fisher metric. Additional assumptions may be
required for such an identification.
\end{remark}
\section{Fisher Information Metric}
Let $\Theta$ be a smooth statistical model with probability distributions
$p(x;\theta).$
The Fisher information metric is defined by
\[
g_F(\theta)
=
\mathbb E_\theta
\!\left[
\nabla_\theta\log p(X;\theta)
\nabla_\theta\log p(X;\theta)^T
\right].
\]
It is a positive-semidefinite symmetric tensor and, under regularity
assumptions, defines a Riemannian metric on the statistical manifold.
\section{Comparison of the Two Geometries}
The pullback metric constructed in Chapter 7 is
$g = F^{*}g_E,$
where
$F = Q\circ L$
is the canonical logarithmic coordinate map.
Unlike the Fisher metric, which is obtained from expectations of score
functions, the pullback metric is induced by the differential of a
globally defined logarithmic embedding modulo gauge equivalence.
Thus the two constructions originate from different geometric
principles.
\begin{proposition}[Shared Properties]
\label{ch:prop:10-1}
Both the Fisher metric and the pullback metric are symmetric
positive-definite tensor fields whenever they are defined on regular
statistical models.
\end{proposition}
\begin{proof}
The Fisher metric is symmetric by construction and positive definite
under the usual regularity assumptions.
The pullback metric is symmetric and positive definite by Theorem 7.6.
\end{proof}

\begin{proposition}[Coordinate Invariance]
\label{ch:prop:10-2}
Both metrics are invariant under smooth reparameterizations of the
underlying manifold.
\end{proposition}
\begin{proof}
The Fisher metric is invariant under smooth coordinate transformations
by classical information geometry.
The pullback metric is defined as the pullback of the Euclidean metric
by a diffeomorphism and is therefore intrinsically
coordinate-independent. 
\end{proof}
\section{Structural Differences}
The Fisher metric depends on
\begin{itemize}
\item the statistical model,
\item the likelihood function,
\item expectations with respect to the model.
\end{itemize}
The pullback metric developed in this work depends on
\begin{itemize}
\item posterior distributions,
\item logarithmic coordinates,
\item quotient geometry induced by additive gauge symmetry.
\end{itemize}
Consequently, the pullback metric remains meaningful even in situations
where no differentiable statistical parameterization is explicitly
available.
\section{Compatibility}
\begin{theorem}[Intrinsic Definition]
\label{ch:thm:10-3}
The pullback metric is intrinsically defined on the posterior manifold
independently of any statistical parameterization.
\end{theorem}
\begin{proof}
By Chapter 6, the canonical logarithmic coordinate map
is a global diffeomorphism onto its image.
The pullback metric is therefore defined entirely through
which depends only on the posterior manifold and the Euclidean structure
of the target space.
No external parameterization is required. 
\end{proof}

\begin{corollary}[Generality]
\label{ch:cor:10-4}
The differential geometry induced by Weighted Incidence Structures
extends beyond the classical setting of parameterized statistical
models.
\end{corollary}
\begin{proof}
Immediate from Theorem~\ref{ch:thm:10-3}. 
\end{proof}
\section{Discussion}
The Fisher metric and the pullback metric share several fundamental
geometric properties:
\begin{itemize}
\item symmetry;
\item positive definiteness;
\item coordinate invariance;
\item smooth dependence on the underlying manifold.
\end{itemize}
However, they arise from distinct constructions.
The Fisher metric is expectation-based and model-dependent.
The pullback metric introduced in this appendix is induced by the
logarithmic geometry of posterior distributions modulo gauge
transformations.
This distinction suggests that the present framework should be viewed as
complementary to classical Information Geometry rather than as a
replacement for it.
Future work may identify conditions under which the two metrics coincide
or become naturally isometric.
\section{Outlook}
The next chapter investigates entropy-based quantities and derives
bounds expressed in terms of the pullback geometry developed throughout
this appendix.
